\chapter[1993 --- Roberts et al.]{1993: Richard J. Roberts, Phillip A. Sharp}
\label{ch:1993_Roberts}

\section*{Main Contribution}

\textit{Discovered that genes are split (containing introns and exons) and that RNA is spliced, revealing a fundamental mechanism of gene expression in higher organisms.}\cite{nobel-1993,lecture-1993}

\section{Official Citation}

for their discoveries of split genes and RNA splicing

\section{Background and Historical Context}

The central dogma of molecular biology, as articulated by Francis Crick in 1958, holds that genetic information flows from DNA to RNA to protein. In its simplest formulation, the gene is a contiguous stretch of DNA that encodes a continuous mRNA molecule, which in turn is translated directly into a protein. This model, elegant in its simplicity, dominated molecular biology for two decades. It was shattered in 1977, when Richard J. Roberts and Phillip A. Sharp independently discovered that genes in higher organisms are not continuous but are instead interrupted by non-coding sequences called introns, and that the mature mRNA is assembled by splicing together the coding sequences---the exons. This discovery of ``split genes'' and RNA splicing fundamentally changed our understanding of gene structure and expression and has had far-reaching implications for molecular biology and medicine.

Richard John Roberts was born on September 6, 1943, in Derby, England. He studied chemistry at the University of Sheffield, earning his B.Sc.\ in 1965 and his Ph.D.\ in 1968 under the supervision of David M. Brown, working on the chemistry of nucleic acids. Roberts became interested in the molecular biology of gene expression and moved to Cold Spring Harbor Laboratory in New York in 1972, where he worked in the laboratory of James Watson. It was at Cold Spring Harbor that Roberts made his discovery of split genes.

Phillip Allen Sharp was born on March 26, 1944, in Falmouth, Kentucky. He studied chemistry at Union College in Schenectady, New York, and then earned his Ph.D.\ in biochemistry at the University of Illinois at Urbana--Champaign under the supervision of Bob Ley. Sharp joined the Massachusetts Institute of Technology (MIT) in 1974, where he established a laboratory studying the molecular biology of adenovirus---a DNA virus that infects human cells. It was through his work on adenovirus that Sharp made his independent discovery of split genes.

The intellectual context of their discoveries was the intense effort to understand how genes are expressed in eukaryotic cells. By the mid-1970s, molecular biologists had developed powerful tools for studying gene structure, including restriction enzymes, gel electrophoresis, DNA sequencing, and in vitro transcription. The adenovirus genome, which was being intensively studied as a model system for eukaryotic gene expression, provided an ideal experimental platform. Adenovirus has a relatively small genome (approximately 36 kilobases), is easy to grow in culture, and uses the host cell's transcription machinery. Studies of adenovirus gene expression were expected to reveal the basic principles of eukaryotic transcription.

\section{The Discovery/Contribution}

\subsection{Richard J. Roberts and Split Genes}

Roberts was studying the late region of the adenovirus-2 genome---the portion of the genome that encodes the structural proteins produced during the late phase of viral infection. Using electron microscopy, he examined hybrids formed between viral DNA and late viral mRNA. If the mRNA were a faithful copy of the DNA, the DNA--RNA hybrid would be perfectly complementary and the two molecules would anneal completely. Instead, Roberts observed that the mRNA contained sequences not present in the contiguous DNA template. When DNA--RNA hybrids were visualized under the electron microscope, loops of single-stranded DNA were seen bulging out from the hybrid duplex---regions of the DNA that had no complementary sequences in the mRNA.

These ``DNA loops'' or ``R-loops'' demonstrated that the gene was discontinuous: the DNA contained intervening sequences (introns) that were present in the initial RNA transcript but were absent from the mature mRNA. The mRNA was assembled by removing these intervening sequences and joining together the flanking coding sequences (exons). Roberts published this finding in a landmark paper in \textit{Cell} in 1977.

\subsection{Phillip A. Sharp and the Confirmation of Splicing}

Independently and almost simultaneously, Sharp and his colleagues at MIT were conducting similar experiments with adenovirus. Using a different approach---analyzing the hybridization patterns of viral RNA to restriction fragments of viral DNA---Sharp observed anomalous hybridization results that could only be explained if the mature mRNA contained sequences derived from non-contiguous regions of the viral genome. Sharp's group also used electron microscopy to visualize DNA--RNA hybrids and observed the characteristic R-loop structures.

Sharp's paper was published in the same issue of \textit{Cell} as Roberts's, and the two discoveries provided independent and complementary confirmation of the existence of split genes and RNA splicing.

\subsection{The Mechanism of RNA Splicing}

The discovery of split genes immediately raised the question: how are introns removed and exons joined? In the years following the initial discovery, researchers identified the molecular machinery responsible for splicing. The spliceosome---a large ribonucleoprotein complex composed of five small nuclear ribonucleoproteins (snRNPs, pronounced ``snurps'') and numerous associated proteins---was identified as the catalytic machine that carries out splicing.

Splicing occurs at specific sequences at the intron--exon boundaries: a conserved GU dinucleotide at the 5' splice site, a conserved AG dinucleotide at the 3' splice site, and a branch point sequence (containing an adenosine residue) located 18--40 nucleotides upstream of the 3' splice site. The spliceosome recognizes these sequences through base pairing between the snRNAs and the pre-mRNA, and catalyzes two transesterification reactions that remove the intron as a lariat (loop) structure and ligate the exons.

A critical consequence of RNA splicing is that a single gene can give rise to multiple different mRNAs---and therefore multiple different proteins---by alternative splicing. In alternative splicing, different combinations of exons are included in the mature mRNA, allowing a single gene to encode a family of related but functionally distinct proteins. This process, which was not anticipated by the simple ``one gene--one protein'' model, vastly expands the coding capacity of the genome. The human genome contains approximately 20,000 protein-coding genes, but alternative splicing allows these genes to produce an estimated 100,000 or more distinct proteins.

\section{Impact on Modern Medicine}

The discovery of split genes and RNA splicing has had significant implications for molecular biology and medicine.

\subsection{Understanding Genetic Disease}

Mutations that disrupt splicing are a major cause of genetic disease. Approximately 15--50\% of all disease-causing mutations affect splicing, either by creating new splice sites, destroying existing ones, or altering splicing regulatory elements. Spinal muscular atrophy (SMA), caused by mutations in the SMN1 gene, is a devastating neuromuscular disease that has been successfully treated by nusinersen (Spinraza), an antisense oligonucleotide that modulates splicing of the paralogous SMN2 gene to increase production of functional SMN protein. This drug, approved in 2016, is a direct clinical application of the understanding of splicing derived from Roberts and Sharp's discovery.

\subsection{Cancer Biology}

Aberrant splicing is increasingly recognized as a hallmark of cancer. Many oncogenic mutations affect splicing factors or splice sites, leading to the production of aberrant protein isoforms that promote cell proliferation, inhibit apoptosis, or enhance metastasis. The discovery of these splicing defects has opened new avenues for cancer diagnosis and therapy.

\subsection{RNA-Based Therapeutics}

The understanding of RNA splicing has also enabled the development of RNA-based therapeutics. Antisense oligonucleotides, small interfering RNAs (siRNAs), and splice-switching oligonucleotides are all therapeutic modalities that depend on a detailed understanding of RNA processing. The approval of multiple RNA-based drugs in recent years---including nusinersen for SMA, eteplirsen for Duchenne muscular dystrophy, and patisiran for hereditary transthyretin amyloidosis---reflects the translation of splicing biology into clinical practice.

\section{Legacy}

Richard J. Roberts continued his work at New England Biolabs, where he has served as chief scientific officer and has made significant contributions to the understanding of restriction enzymes and DNA methylation. He has been a vocal advocate for science education and for the role of basic research in driving medical progress.

Phillip A. Sharp became director of the McGovern Institute for Brain Research at MIT and has been a leader in the fields of RNA biology and gene regulation. He has received numerous honors, including the National Medal of Science and the Albert Lasker Basic Medical Research Award.

Together, Roberts and Sharp revealed a fundamental truth about the architecture of genes that challenged the prevailing model of gene expression. Their discovery of split genes and RNA splicing opened entirely new fields of research---from the spliceosome and alternative splicing to RNA therapeutics and splicing-based diagnostics---that continue to advance our understanding of biology and to improve human health.


\section{Modern Developments}

Roberts and Sharp's split gene discovery has led to modern RNA biology. Understanding of splicing has revealed that ~95% of human genes undergo alternative splicing, creating protein diversity. Understanding of splicing has led to treatments for spinal muscular atrophy (nusinersen, which modulates splicing of SMN2). RNA-based therapeutics, including antisense oligonucleotides and siRNAs, target splicing and RNA processing. Understanding of aberrant splicing in cancer has led to new diagnostic and therapeutic approaches.


\section{Bibliography}

\begin{thebibliography}{9}
\bibitem{nobel-1993}
Nobel Prize Outreach, ``The Nobel Prize in Physiology or Medicine 1993,''\emph{NobelPrize.org}.\\
\url{https://www.nobelprize.org/prizes/medicine/1993/summary/}

\bibitem{lecture-1993}
Richard J. Roberts, ``Nobel Lecture,''\emph{NobelPrize.org}. Winner-authored primary source; downloadable from the official Nobel Prize archive.\\
\url{https://www.nobelprize.org/prizes/medicine/1993/roberts/lecture/}
\end{thebibliography}
